\documentclass[../main.tex]{subfiles}

\begin{document}

\section{Task 2: Vertical Climb and Optimal Burn}\label{sec:task2}

\subsection{Phase 1: Vertical Climb}

Real launchers cannot pitch over right at liftoff, both because of aerodynamic loads and because they have to clear the tower first. Task~2 adds a short vertical-climb phase in front of the burn ($\varphi = \pi/2$, thrust pointing straight up) until the altitude reaches $y_1 = 10^{-4}$ (dimensionally $\approx 620$~m). During this phase the horizontal position and velocity remain zero, and the 3-state sub-problem is
\begin{equation}
    \dot{y} = v_y, \qquad
    \dot{v}_y = \frac{T}{m} - 1, \qquad
    \dot{m} = -Q,
    \label{eq:ode_vert}
\end{equation}
integrated until the event $y = y_1$ is detected via \texttt{ode45} event handling. The end state $(y_1,\, v_{y,1},\, m_1)$ becomes the initial condition for Phase~2.

\subsection{Phase 2: Optimal Burn}

Starting from the state $(x_0, y_0, v_{x,0}, v_{y,0}, m_0) = (0,\, y_1,\, 0,\, v_{y,1},\, m_1)$, Phase~2 is an optimal powered arc with the same BVP structure as Task~1. The shooting unknowns and the residual \eqref{eq:residual1} are unchanged; only the initial conditions of the augmented state differ. Because the orbit injection targets are identical to Task~1, the continuation warm-start from Task~1's solution at $Q = 2$ provides a reliable initial guess.

\subsection{Results}

Figure~\ref{fig:task2} shows the composite trajectory and the angle history for $Q = 2$, $y_f = 0.04$. The red segment is the short vertical-climb phase; the blue arc is the optimal burn. The trajectory is qualitatively similar to Task~1 except for the initial vertical stub. During the climb both $\varphi$ and $\psi$ are fixed at $90^\circ$; at the transition the thrust begins to rotate toward the horizontal following the bilinear-tangent law, while the velocity angle $\psi$ lags momentarily before catching up as horizontal speed builds.

\begin{figure}[H]
    \centering
    \begin{subfigure}[b]{0.68\textwidth}
        \centering
        \includegraphics[width=\textwidth]{task2_task_2_trajectory}
        \caption{Composite trajectory: vertical climb (red), optimal burn (blue).}
        \label{fig:task2_traj}
    \end{subfigure}
    \\[2ex]
    \begin{subfigure}[b]{0.68\textwidth}
        \centering
        \includegraphics[width=\textwidth]{task2_task_2_angles}
        \caption{Thrust angle $\varphi$ and velocity angle $\psi$ during the burn.}
        \label{fig:task2_angles}
    \end{subfigure}
    \caption{Task 2: trajectory and control history for $Q = 2$, $y_f = 0.04$.}
    \label{fig:task2}
\end{figure}

The final mass in Task~2 is marginally lower than in Task~1: the vertical climb consumes propellant while building only vertical velocity, contributing to gravity losses without advancing the vehicle toward the horizontal injection speed. Nevertheless, the prepended climb is a standard launch practice, and the payload penalty is modest for the short $y_1$ value used here. The converged simulation results are collected in Table~\ref{tab:task2}.

\begin{table}[H]
\centering
\caption{Task 2 simulation results ($Q = 2$, $y_f = 0.04$). All quantities are non-dimensional; the Task~1 final mass at the same $Q$ is reported for reference.}
\label{tab:task2}
\begin{tabular}{llc}
\toprule
Quantity                       & Symbol            & Value \\
\midrule
Vertical-climb duration        & $t_1$             & $0.0298$ \\
Powered-burn duration          & $t_\mathrm{burn}$ & $0.4142$ \\
Total flight time              & $t_f$             & $0.4440$ \\
Final mass                     & $m_f$             & $0.1119$ \\
Delivered payload              & $m_u$  & $0.0231$ \\
\midrule
Final mass --- Task 1 (no climb) & $m_f^{(1)}$     & $0.1128$ \\
\bottomrule
\end{tabular}
\end{table}

\end{document}
